% Problem record op_3b3bfcda365a83a9. This TeX file is derived from
% database/problems_json/op_3b3bfcda365a83a9.json by scripts/sync-tex.mjs; edit the JSON record
% and rerun the script rather than editing this file.
\section{Linear-size exact quantum Fourier transform}
\label{sec:op_3b3bfcda365a83a9}

\paragraph{Problem.}
Can the exact quantum Fourier transform on $n$ qubits be implemented with $O(n)$ one- and two-qubit gates?

Define

\begin{equation}
F_{2^n}|x\rangle=2^{-n/2}\sum_{y=0}^{2^n-1}e^{2\pi ixy/2^n}|y\rangle.
\label{eq:3b3b-1}
\end{equation}

Let $C_{\mathrm{exact}}(n)$ be the minimum size of a uniform circuit that implements the unitary in Eq.~\eqref{eq:3b3b-1} exactly on arbitrary inputs. Gates may be arbitrary efficiently specified one- or two-qubit unitaries; clean ancillas may be used but must be restored, and all gates on them count. Is

\begin{equation}
C_{\mathrm{exact}}(n)=O(n)?
\label{eq:3b3b-2}
\end{equation}

The target in Eq.~\eqref{eq:3b3b-2} concerns exact coherent implementation, not approximate QFT or sampling only.

\subsection*{Status}

Unsolved

\subsection*{Source}

Aaronson explicitly asks the linear-size exact-QFT question in the cited author-written discussion \sourcecite{ref:3b3b-aaronson25}{Aaronson25}. Cleve and Watrous and Kahanamoku--Meyer and Yao supply the formal circuit results \sourcecite{ref:3b3b-cleve00}{Cleve00}\sourcecite{ref:3b3b-kmy24}{KMY24}.

\subsection*{Progress}

\begin{itemize}
  \item Cleve and Watrous obtained an exact QFT of size $O(n\log^2 n\log\log n)$ in 2000, together with efficient approximate parallel constructions. Their exact construction uses recursive Fourier transforms and fast multiplication. Thus a subquadratic exact QFT has long been known. \sourcecite{ref:3b3b-cleve00}{Cleve00}

  \item Kahanamoku--Meyer and Yao subsequently constructed zero-ancilla exact QFT circuits of size $O\!\left(n^{\log_k(2k-1)}\right)$ for every fixed integer $k\geq2$. Since $\log_k(2k-1)$ approaches $1$ as $k$ grows, this gives size $O(n^{1+\varepsilon})$ for every fixed $\varepsilon>0$. It does not give a single uniform $O(n)$ construction. \sourcecite{ref:3b3b-kmy24}{KMY24}

  \item Aaronson explicitly raised the $O(n)$ exact-size question in January 2025. The construction above improves the upper bound for zero-ancilla exact QFT, but the linear target remains unresolved. \sourcecite{ref:3b3b-aaronson25}{Aaronson25}

  \item Shah withdrew \emph{A Faster Quantum Fourier Transform} in February 2025 because of lack of novelty. The withdrawal does not bear on whether linear exact size is possible. \sourcecite{ref:3b3b-shah25}{Shah25}

  \item Fault-tolerant and architectural results use different metrics. Nam, Su, and Maslov obtained an approximate QFT with $O(n\log n)$ $T$ gates, while Lopes studies QFT execution using phase-gradient resources, surface codes, and resource routing. Neither establishes linear exact coherent circuit size in the model above. \sourcecite{ref:3b3b-nam20}{Nam20}\sourcecite{ref:3b3b-lopes26}{Lopes26}
\end{itemize}

\subsection*{References}

\begin{enumerate}[label={},leftmargin=4.8em,labelsep=0.5em]
  \footnotesize
  \item[\textup{[Cleve00]}]\label{ref:3b3b-cleve00}
  Richard Cleve and John Watrous. \emph{Fast parallel circuits for the quantum Fourier transform}. FOCS 2000; \href{https://arxiv.org/abs/quant-ph/0006004}{arXiv:quant-ph/0006004}. See Theorem 2 and the exact-construction recurrence.

  \item[\textup{[KMY24]}]\label{ref:3b3b-kmy24}
  Gregory D. Kahanamoku--Meyer and Norman Y. Yao, \emph{Fast quantum integer multiplication with zero ancillas}, arXiv:2403.18006v4 (14 November 2024). \href{https://arxiv.org/abs/2403.18006}{arXiv:2403.18006}. See the exact QFT construction and its $O(n^{1+\varepsilon})$ consequence.

  \item[\textup{[Aaronson25]}]\label{ref:3b3b-aaronson25}
  Scott Aaronson. \href{https://scottaaronson.blog/?p=8593}{Author-written QFT discussion dated January 23, 2025}, with subsequent corrections and comments by Richard Cleve. Used for the explicit open question and corrected historical attribution, not as a replacement for \sourcecite{ref:3b3b-cleve00}{Cleve00}.

  \item[\textup{[Shah25]}]\label{ref:3b3b-shah25}
  Ronit Shah. \emph{A Faster Quantum Fourier Transform}. \href{https://arxiv.org/abs/2501.12414}{arXiv:2501.12414}; withdrawn February 10, 2025 for lack of novelty.

  \item[\textup{[Nam20]}]\label{ref:3b3b-nam20}
  Yunseong Nam, Yuan Su, and Dmitri Maslov. \emph{Approximate Quantum Fourier Transform with $O(n\log(n))$ T gates}. \href{https://www.nature.com/articles/s41534-020-0257-5}{npj Quantum Information 6, 26 (2020)}.

  \item[\textup{[Lopes26]}]\label{ref:3b3b-lopes26}
  Pedro L. S. Lopes. \emph{Towards Deploying Optimistic Quantum Fourier Transforms: An Architecture-Algorithm Co-Design Study}. \href{https://arxiv.org/abs/2605.15297}{arXiv:2605.15297}, May 14, 2026, preprint.
\end{enumerate}

\subsection*{Comment}

This is a literature-explicit circuit-complexity question. The known exact upper bound is $O(n^{1+\varepsilon})$ for every fixed $\varepsilon>0$, whereas no superlinear lower bound is known in the stated arbitrary-gate model. Exact coherent QFT computation is stronger than what bounded-error factoring needs. A resolution must therefore either improve the family of $O(n^{1+\varepsilon})$ upper bounds to $O(n)$ or prove that some superlinear growth is unavoidable.

\subsection*{Field}

Quantum algorithm

\subsection*{Topic}

Quantum circuit complexity

\subsection*{ID}

\texttt{op\_3b3bfcda365a83a9}
